\documentclass[12pt]{ximera}
\title{Homework 4: Weighted Voting Systems and Banzhaf Power}
\begin{document}
\begin{abstract}
Sections 2.1--2.2: equivalent weighted voting systems and dummies, the Banzhaf power of a
law firm's senior partner, and winning coalition types for a city council with a
mayoral veto.
\end{abstract}
\maketitle

\section*{Sections 2.1--2.2: Weighted Voting and Banzhaf Power}

Recall the quota restrictions: in $[q: w_1, w_2, \ldots, w_n]$ the quota $q$ must be
\emph{more than half} of the total weight and \emph{no more than} the total weight. A
player whose weight alone meets the quota is a \emph{dictator}.

\begin{problem}
Two weighted voting systems are \emph{equivalent} if and only if they have the same
number of players and the same winning coalitions.

Consider the weighted voting system $[8:5,3,2,1]$, with players A, B, C, and D.

\begin{prompt}
\textbf{(a)} Is this weighted voting system equivalent to $[2:1,1,0,0]$?

\begin{multipleChoice}
\choice{Yes}
\choice[correct]{No}
\end{multipleChoice}

\begin{hint}
In $[2:1,1,0,0]$, a coalition wins exactly when it contains both A and B. Can you find
a winning coalition in $[8:5,3,2,1]$ that does \emph{not} contain B?
\end{hint}

\begin{feedback}
No. In $[2:1,1,0,0]$ only A and B have any weight, so a coalition wins precisely when
it contains \emph{both} A and B.

In $[8:5,3,2,1]$ the winning coalitions (those with weight at least $8$) are
\[
\set{A,B}=8,\quad \set{A,B,C}=10,\quad \set{A,B,D}=9,\quad \set{A,B,C,D}=11,\quad
\set{A,C,D}=8.
\]
The last one is the counterexample: $\set{A,C,D}$ has weight $5+2+1=8$ and wins
\emph{without} B. In $[2:1,1,0,0]$ that same coalition has weight $1$ and loses.

The two systems have different winning coalitions, so they are not equivalent.
\end{feedback}
\end{prompt}

\begin{prompt}
\textbf{(b)} Are there any dummies in $[8:5,3,2,1]$?

\begin{multipleChoice}
\choice{Yes --- D is a dummy}
\choice{Yes --- C and D are both dummies}
\choice[correct]{No --- every player is critical in at least one winning coalition}
\end{multipleChoice}

\begin{feedback}
No. A dummy is a player who is never critical in any winning coalition.

The coalition $\set{A,C,D}$ settles it for both of the small players. Its weight is
exactly $8$, so removing either C or D drops it to $6$ or $7$ and it loses. Both C and
D are therefore critical in that coalition, and neither is a dummy.

(A and B are clearly not dummies either --- both are critical in $\set{A,B}$.)
\end{feedback}
\end{prompt}
\end{problem}

\begin{problem}
A law firm is run by four partners --- A, B, C, and D --- where A is the senior
partner. Each partner has one vote. Decisions are typically made by majority rule, but they can
also be made by the senior partner with any one other partner's support.

\begin{prompt}
\textbf{(a)} How many winning coalitions are there in this system?

$\answer{8}$

\begin{feedback}
A majority of four votes is three, so the winning coalitions are every group of three or
more, plus the senior partner with any one other:
\[
\set{A,B},\ \set{A,C},\ \set{A,D},\ \set{A,B,C},\ \set{A,B,D},\ \set{A,C,D},\
\set{B,C,D},\ \set{A,B,C,D},
\]
which is $8$ coalitions. Marking the critical players in each --- those whose removal
would make the coalition lose:
\[
\begin{array}{ll}
\set{A,B}=3: & \text{both A and B critical}\\
\set{A,C}=3: & \text{both A and C critical}\\
\set{A,D}=3: & \text{both A and D critical}\\
\set{A,B,C}=4: & \text{only A critical}\\
\set{A,B,D}=4: & \text{only A critical}\\
\set{A,C,D}=4: & \text{only A critical}\\
\set{B,C,D}=3: & \text{all three of B, C, D critical}\\
\set{A,B,C,D}=5: & \text{nobody critical}
\end{array}
\]
\end{feedback}
\end{prompt}

\begin{prompt}
\textbf{(b)} Determine the Banzhaf Power Distribution of this law firm. Give each
share as a percentage, rounded to one decimal place.

A: $\answer[tolerance=0.05]{50}\%$ \qquad
B: $\answer[tolerance=0.05]{16.7}\%$ \qquad
C: $\answer[tolerance=0.05]{16.7}\%$ \qquad
D: $\answer[tolerance=0.05]{16.7}\%$

\begin{hint}
Count how many times each player is critical across all winning coalitions, then
divide by the grand total of critical counts.
\end{hint}

\begin{feedback}
From the list in part (b), count each player's critical appearances:
\[
\text{A}: 6, \qquad \text{B}: 2, \qquad \text{C}: 2, \qquad \text{D}: 2,
\]
for a total critical count of $12$. The Banzhaf power distribution is therefore
\[
\text{A}: \frac{6}{12}=50\%, \qquad
\text{B}=\text{C}=\text{D}: \frac{2}{12}=16.7\%.
\]
(These total $100\%$. \checkmark)

The senior partner holds exactly half the power with only $40\%$ of the raw votes ---
Banzhaf power is not proportional to weight.
\end{feedback}
\end{prompt}
\end{problem}

\begin{problem}
A city council has seven council members, $C_1,\ldots,C_7$, with one vote each, and a
mayor ($M$) with no vote but a veto, which a two-thirds majority of the council can
overrule. A simple majority of the council is needed to pass a motion. Use $C$ for any
one council member, and include $M$ in a coalition when the mayor does not intend to
veto.

\begin{prompt}
\textbf{(a)} How many council members does a coalition need to pass a motion
\emph{with} mayoral support?

$\answer{4}$

\begin{feedback}
A simple majority of $7$ council members means more than $3.5$, so at least $4$.
\end{feedback}
\end{prompt}

\begin{prompt}
\textbf{(b)} How many council members does a coalition need to pass a motion
\emph{without} mayoral support?

$\answer{5}$

\begin{hint}
Without mayoral support the motion is vetoed, so the council must override. Two-thirds
of $7$ is not a whole number --- round up.
\end{hint}

\begin{feedback}
Overriding the veto takes a two-thirds majority of the $7$ council members, and
\[ \tfrac23 \cdot 7 = 4.67, \]
so at least $5$ council members are required.
\end{feedback}
\end{prompt}

\begin{prompt}
\textbf{(c)} Below is every coalition type. Identify which are winning, and how many
players are critical in each.

For coalitions \emph{with} mayoral support:

$\set{M, C, C, C, C}$ (4 council members, $35$ distinct coalitions) --- is it winning?
\begin{multipleChoice}
\choice[correct]{Yes}
\choice{No}
\end{multipleChoice}
Number of critical players in each such coalition: $\answer{5}$

\smallskip
$\set{M, C, C, C, C, C}$ (5 council members, $21$ coalitions) --- number of critical
players: $\answer{0}$

\smallskip
For coalitions \emph{without} mayoral support:

$\set{C, C, C, C, C}$ (5 council members, $21$ coalitions) --- number of critical
players: $\answer{5}$

\smallskip
$\set{C,C,C,C,C,C}$ (6 council members, $7$ coalitions) --- number of critical
players: $\answer{0}$

\begin{hint}
A player is critical when removing them turns a winning coalition into a losing one.
Careful with $\set{M,C,C,C,C,C}$: if the mayor walks away, the five remaining council
members can still override his veto.
\end{hint}

\begin{feedback}
The winning types are: with the mayor, $4$ or more council members; without the mayor,
$5$ or more.

\textbf{$\set{M,C,C,C,C}$ --- winning, $5$ critical players.} Remove $M$ and the four
remaining council members fall short of the $5$ needed to override, so $M$ is
critical. Remove any one $C$ and only $3$ council members remain with mayoral support,
short of $4$ --- so all four council members are critical too. That is $1+4 = 5$
critical players.

\textbf{$\set{M,C,C,C,C,C}$ --- winning, but nobody is critical.} Remove $M$ and the
$5$ council members override the veto and still win. Remove one $C$ and you are left
with $\set{M,C,C,C,C}$, which also still wins. The same holds for every larger
coalition with the mayor.

\textbf{$\set{C,C,C,C,C}$ --- winning, all $5$ critical.} Removing any member leaves
$4$, which cannot override the veto.

\textbf{$\set{C,C,C,C,C,C}$ --- winning, nobody critical.} Removing one still leaves
$5$, enough to override.

Every other type is either losing (too few council members) or winning with no
critical players.
\end{feedback}
\end{prompt}

It turns out that both the mayor and each individual council member hold exactly
$12.5\%$ of the Banzhaf power in this system --- the mayor's veto is worth exactly as much
as a council vote.
\end{problem}

\end{document}
